\documentclass[10pt,a4paper]{amsart}
\usepackage[margin=19mm]{geometry}
\usepackage{amsmath,amssymb,mathtools}
\usepackage[scr=rsfs,cal=euler]{mathalfa}
\usepackage{xfrac}
\usepackage[unicode,hidelinks,bookmarks=false]{hyperref}
\newcommand{\ie}{i.e.\ }
\newcommand{\cf}{cf.\ }
\newcommand{\resp}{respectively\ }
\newcommand{\Subsection}{\subsection}
\providecommand{\nobreakdash}{\nobreak\hbox{-}}
\setlength{\emergencystretch}{2em}
\multlinegap0pt
\usepackage{bm}
\usepackage{bbm}
\usepackage{dsfont}
\usepackage{xspace}
\usepackage{cancel}
\usepackage{mathtools}
\usepackage[shortlabels]{enumitem}
\usepackage{amsthm}
\makeatletter
\@ifundefined{theorem}{\newtheorem{theorem}{Theorem}}{}
\@ifundefined{lemma}{\newtheorem{lemma}[theorem]{Lemma}}{}
\@ifundefined{proposition}{\newtheorem{proposition}[theorem]{Proposition}}{}
\makeatother
\newcommand{\NH}{\mathcal N}
\newcommand{\V}{\mathcal V}
\newcommand{\disc}{\operatorname{disc}}
\title[Bang--bang representation of $3\times 3$ embeddable stochast]{Bang--bang representation of $3\times 3$ embeddable stochastic matrices}
\author{Leonel Robert}
\date{Source: arXiv:2608.08242v1 (CC BY 4.0)}
\begin{document}
\maketitle

\noindent\emph{Source and licence.} Bang--bang representation of $3\times 3$ embeddable stochastic matrices. Version arXiv:2608.08242v1 (CC BY 4.0), distributed under the Creative Commons Attribution 4.0 International licence, as recorded in the arXiv licence metadata for this exact version and shown on the abstract page https://arxiv.org/abs/2608.08242v1. Full rights notice: licenses/arxiv-2608.08242-CC-BY-4.0.txt.\par\smallskip

\noindent\textit{Editorial scope.} Counted unit: the Lemma whose source label is \texttt{lem:R6-elimination} in the article (source0.tex lines 1197--1206, source label \texttt{lem:R6-elimination}), delivered together with the article's own complete original proof of it (source0.tex lines 1208--1421, one \texttt{proof} environment of 214 lines). No mathematics is written, repaired or re-derived: statement and proof are the article's own text line for line. Required context retained uncounted: prop:Rn-relatively-closed (lines 460--463); prop:reduced-Frydman-test (lines 773--786); prop:last-move-quadratic-criterion (lines 875--880); bysymmetry (lines 1175--1179). Every definition, display and label of those lines is retained unchanged. Omitted from this package and not redistributed: 1--459, 464--772, 787--874, 881--1174, 1180--1196, 1207--1207, 1422--2329. Nothing inside a retained excerpt is omitted or altered: every retained body is delivered exactly as the article's lines are written, so no omission receipt of any kind accompanies this package. Citations: the retained excerpts carry no citation command at all, so no bibliography entry of the article is transcribed and no reference is redistributed. Reference adaptation: none was needed and none was made; every reference inside the delivered unit resolves inside it, so no unresolved reference or citation is produced. Printed slips kept literal: no printed word, symbol, hypothesis or number of the counted statement or of its proof was altered. Layout and class: the article's own class is \texttt{amsart} with packages including geometry, fontenc, lmodern, microtype, amsmath, amssymb, amsthm, aliascnt, mathtools, fancyhdr, hyperref, cleveref, enumitem, tikz; the wrapper is the lane's pinned \texttt{amsart} editorial wrapper at 10pt on A4, which restates the theorem-like environments the retained text needs and adds the article's own macro definitions that the retained text needs (\texttt{\textbackslash NH}, \texttt{\textbackslash V}, \texttt{\textbackslash disc}), with extra wrapper packages (bm, bbm, dsfont, xspace, cancel, mathtools, enumitem). The delivered numbering is the unit's own. Assets: no retained span contains a figure environment, a graphics-inclusion command, a PDF-page inclusion, a table or a TikZ picture; no image file is copied into this package, the package declares no asset dependency and no image file is redistributed with it. Licence: version arXiv:2608.08242v1 of this article is distributed under the Creative Commons Attribution 4.0 International licence, as recorded in the arXiv licence metadata for this exact version and shown on the abstract page https://arxiv.org/abs/2608.08242v1. A full rights notice is delivered at licenses/arxiv-2608.08242-CC-BY-4.0.txt. Characters: every retained excerpt is pure ASCII, so the delivered engine, XeLaTeX with its default Latin Modern text font, reports no missing character.
\begin{proposition}
\label{prop:Rn-relatively-closed}
For every $n$, the set $R_n$ is relatively closed in $\mathcal M_+$.
\end{proposition}

\begin{proposition}[Reduced Frydman test]
\label{prop:reduced-Frydman-test}
Assume that $A>0$ and $A\notin\V$.  Fix a permutation $(i,j,k)$ and
write $C=C_{ij}$.  Then $q_{ij}^A$ has a root in $[0,C]$ if and only if
\[
T_{jj}<0,
\qquad
\beta_{ij}>0,
\qquad
2\alpha_{ij}C+\beta_{ij}<0,
\qquad
T_{jk}<0,\qquad \widetilde H_{ij}\leq0.
\]
\end{proposition}

\begin{proposition}
	\label{prop:last-move-quadratic-criterion}
	Suppose that \(A>0\) and \(A\notin\V\). Then \(A\in\NH\) if and only
	if, for some ordered pair \((i,j)\), the quadratic \(G_{ij}^A\) has a
	root in \([0,S_{ij}]\).
\end{proposition}

\begin{equation}\label{bysymmetry}
U(x,r,y)\in R_n
\quad\Longleftrightarrow\quad
U(y,r,x)\in R_n.
\end{equation}

\begin{lemma}[Six-move elimination]
\label{lem:R6-elimination}
Suppose that $U\notin\V$.  Then $U\in R_6$ in each of the following
regions:
\begin{enumerate}[(1)]
\item $1<r\leq2$ and $\min(x,y)\leq r$;
\item $r\geq3$ and $\max(x,y)\geq r$;
\item $2\leq r\leq3$ and $\min(x,y)\geq r$.
\end{enumerate}
\end{lemma}

\begin{proof}
Since $U\notin\V$, every $V_{ij}(U)$ is negative.

\smallskip
\noindent
\emph{Case (1).}
Assume first that $1<r<2$ and $\min(x,y)<r$. By the symmetry \eqref{bysymmetry}, we may assume that $y<r$. We apply the
last-move quadratic test from
Proposition~\ref{prop:last-move-quadratic-criterion} to
\[
G_{12}^U(t)=V_{21}(X_{12}(-t)U)=\alpha t^2+\beta t+\gamma,
\]
whose coefficients are
\[
\alpha=(x+1)E,
\qquad
\beta=-2(x+1)(yE-x),
\qquad
\gamma=V_{21}(U).
\]
Since 
\[
E=T_{11}(U)=x(r-3)-(r-1)<0,
\]
we have  $\alpha<0$ and $\beta>0$. Since $U\notin \V$, we also have $\gamma=V_{21}(U)<0$.
A direct calculation shows that
\[
\disc(G_{12}^U)
=4x^2(r-2)(x+1)(rx-r-2x).
\]
Since $r-2<0$ and
\[
rx-r-2x=x(r-2)-r<0,
\]
the discriminant $\disc(G_{12}^U)$ is positive.

The three candidates for the admissible endpoint are
\[
\frac{r+2y}{2},
\qquad
y+1,
\qquad
\frac{yx+y+r+x}{x+1}.
\]
Under the present hypotheses, the latter two are strictly larger than the first.
Hence the admissible endpoint is
\[
\bar c=\frac{r+2y}{2}.
\]
The vertex of $G_{12}^U$ is
\[
\tau=\frac{yE-x}{E},
\]
and
\[
\frac{r+2y}{2}-\tau
=\frac{(r-1)(rx-r-2x)}{2E}>0.
\]
Thus $G_{12}^U$ has a root in its admissible interval. By Proposition \ref{prop:last-move-quadratic-criterion}, $U\in R_6$. Proposition~\ref{prop:Rn-relatively-closed} gives the
boundary cases.


\smallskip
\noindent
\emph{Case (2).}
Assume $r>3$ and, by symmetry, $y>r$. We apply the reduced Frydman test
in Proposition~\ref{prop:reduced-Frydman-test} with ordered pair
$(2,1)$. Write
\[
q_{21}^U(t)=\alpha t^2+\beta t+\gamma.
\]
Then
\[
\alpha=(x+1)(yx+y+r+x)E,
\qquad
\beta=-(x+1)\Psi,
\qquad
\gamma=V_{12}(U),
\]
where
\[
\Psi=2V_{12}(U)+(r-2)(-2y+r-4).
\]
Since $V_{12}(U)<0$, and $y>r>3$, we clearly have $\Psi<0$.
Since
\[
V_{12}(U)-E=y(r-3)(x+1)+2r-5>0,
\]
we also have $E<V_{12}(U)<0$. Hence $\alpha<0$, $\beta>0$, and $\gamma<0$.

The complementary minor is $T_{13}(U)=2-r<0$, and
\[
\begin{aligned}
-\widetilde H_{21}(U)={}&\bigl((r-2)(2y+r)^2-8y-8\bigr)x\\
&+(r-2)(2y+r)^2-8y-8r.
\end{aligned}
\]
The quantity
\[
(r-2)(2y+r)^2-8y
\]
increases with $y$ and is larger than
$8r$ at $y=r$.  Thus $\widetilde H_{21}(U)<0$.

The admissible endpoint is
\[
C_{21}=\frac{2}{r+2x}.
\]
The other candidates are
\[
\frac{y+1}{yx+y+r+x}
\qquad\text{and}\qquad
\frac{1}{x+1},
\]
and direct subtraction gives the required ordering. If
$\tau=-\beta/(2\alpha)$, then
\[
\frac2{r+2x}-\tau
=-\frac{(r-2)(2yrx-2yx+2y+r^2+2r+2x)}
{2(r+2x)(yx+y+r+x)E}>0.
\]
The reduced Frydman test gives $U\in R_6$.
Proposition~\ref{prop:Rn-relatively-closed} gives the boundary cases.

\smallskip
\noindent
\emph{Case (3).}
Assume $2<r<3$ and, by symmetry, $r<y\leq x$. We apply the reduced
Frydman test with ordered pair $(3,2)$. For
\[
q_{32}^U(t)=\alpha t^2+\beta t+\gamma,
\]
one has
\[
\alpha=2(r+1)(y(r-3)-r+1),
\qquad
\beta=-(r+1)\Phi,
\qquad
\gamma=V_{23}(U),
\]
where
\[
\Phi=V_{12}(U)-2\bigl((r-2)(x-y)+2(r-2)+2\bigr)<0.
\]
Thus $\alpha<0$, $\beta>0$, and $\gamma<0$.


Put
\[
\Gamma=T_{21}(U)
=yrx-yr-3yx+y-rx+r+x.
\]
Note that $\Gamma$ is affine in each of $x$ and $y$ separately, and that, for fixed \(2<r<3\), we have
\[
\frac{\partial \Gamma}{\partial x}
=yr-3y-r+1
=y(r-3)-r+1<0
\]
and
\[
\frac{\partial \Gamma}{\partial y}
=rx-r-3x+1
=x(r-3)-r+1<0.
\]
Thus \(\Gamma\) is strictly decreasing in
both \(y\) and \(x\). Since \(r<y\leq x\), it follows that
\[
\Gamma<\Gamma|_{y=x=r}
=r^3-5r^2+3r.
\]
Moreover, on the range $2< r \leq 3$ we have $r^3-5r^2+3r<-6$.
Therefore,
$\Gamma<-6$. Consequently, 
\[
T_{21}(U)=\Gamma<0,
\]
and, since \(r+1>3\),
\[
\widetilde H_{32}(U)
=(r+1)\Gamma+16
<-6(r+1)+16<0.
\]

The admissible endpoint is the smaller of
\[
\frac{yx+y+r+x}{2y+r}
\qquad\text{and}\qquad
\frac{x+1}{2}.
\]
The third candidate is larger
than both, since
\[
\frac{r+2x}{r+1}-\frac{yx+y+r+x}{2y+r}
=-\frac{\Gamma}{(2y+r)(r+1)}>0
\]
and
\[
\frac{r+2x}{r+1}-\frac{x+1}{2}
=-\frac{E}{2(r+1)}>0.
\]
For
$\tau=-\beta/(2\alpha)$,
\[
\frac{x+1}{2}-\tau
=\frac{\Gamma+2}{4(y(r-3)-r+1)}>0,
\]
\[
\frac{yx+y+r+x}{2y+r}-\tau
=-\frac{(r-4-2y)\Gamma}
{4(2y+r)(y(r-3)-r+1)}>0.
\]
The reduced Frydman test again gives $U\in R_6$, and
Proposition~\ref{prop:Rn-relatively-closed} gives the boundary cases.
\end{proof}

\end{document}
